\documentclass{article}
\title{Beyond Ability: Cultural and Psychological Factors Shaping Gender Gaps in Mathematics}
\date{}
\begin{document}

\maketitle
 Although women now make up the majority of university students and often perform very well in school, they remain underrepresented in mathematics and STEM fields. The gap tends to appear as students move into more advanced mathematics, especially in problem-solving tasks. These differences are not the result of innate ability but rather cultural expectations, patterns of socialization, and psychological factors that create unequal conditions for women.

Expectancy-value theory and social cognitive theory help explain why this happens. Both highlight how confidence, perceptions of usefulness, and gendered expectations guide students’ decisions about whether to persist in math. Research shows that girls frequently underestimate their mathematical ability, even when their results match or exceed those of boys, and gradually invest their confidence in other subjects. The way self-efficacy develops also differs: women are more influenced by encouragement and role models, while men tend to emphasize mastery experiences.

Stereotype threat adds another layer, increasing anxiety and consuming mental resources, sometimes leading women to adjust their identities in order to cope. These pressures are reinforced by parents, teachers, and the lack of visible female role models, producing cumulative disadvantages over time.

Addressing these barriers requires interventions that strengthen both competence and confidence, create inclusive classrooms, and provide diverse role models. Further research should continue to investigate how stereotypes shape choices, so policies can better promote gender equity in mathematics education and careers.   
\end{document}
